% ECONOMICS_PROBLEM: {"id": "F44-296", "title": "Global financial-cycle transmission", "jel_code": "F44", "source_id": "OP-0333"}
% !TeX program = xelatex
\documentclass[11pt,a4paper]{article}
\usepackage[margin=23mm]{geometry}
\usepackage{fontspec}
\setmainfont{Latin Modern Roman}
\usepackage{amsmath,amssymb,mathtools}
\usepackage{xcolor,enumitem,booktabs,longtable,array,xurl,hyperref}
\definecolor{muted}{HTML}{53636C}
\hypersetup{colorlinks=true,urlcolor=blue,linkcolor=blue}
\setlength{\parindent}{0pt}
\setlength{\parskip}{5pt}
\newcommand{\R}{\mathbb R}
\newcommand{\E}{\mathbb E}
\newcommand{\N}{\mathbb N}
\newcommand{\ind}{\mathbf 1}
\newcommand{\Var}{\operatorname{Var}}
\newcommand{\Cov}{\operatorname{Cov}}
\newcommand{\supp}{\operatorname{supp}}
\DeclareMathOperator*{\argmin}{arg\,min}
\DeclareMathOperator*{\argmax}{arg\,max}
\newcommand{\entrymeta}[3]{{\sffamily\small\color{muted}\textbf{\texttt{#1}}\quad|\quad #2\par #3\par}}
\newcommand{\Problemref}[1]{\texttt{#1}}
\newcommand{\JELref}[2]{\texttt{#2} (JEL \texttt{#1})}
\begin{document}
\section*{Global financial-cycle transmission}
\textbf{Problem ID:} \texttt{F44-296}\quad \textbf{JEL:} \texttt{F44}\par
% BEGIN ECONOMICS STATEMENT
\label{problem:OP-0333}
{\sffamily\small\color{muted}Original subject group: International finance and exchange rates\par}
\label{definitions:problem:30007615}
\textbf{Audit classification:} Published research agenda. The agenda supports global-financial-cycle channels. The safe-asset shock normalization and channel-removal causal definitions are editorial.
\subsection*{Definitions and precise research target}
Fix countries \(c=1,\ldots,C\), intermediary balance sheets, dollar funding exposures and demand for safe dollar claims. Let \(z_t\) be an externally identified global safe-asset-demand shock, normalized to unit variance. Let \(Y_{c,t+h}\) denote a declared outcome such as output, local credit or asset prices. Define total responses \(\tau_c(h)=\partial E[Y_{c,t+h}(do(z_t=u))]/\partial u|_{u=0}\). To separate funding and nonbank channels, specify counterfactual regimes that hold each channel's transmission equation at its pre-shock value while allowing all remaining equations to adjust. The task is to identify total and channel-specific responses, including interactions that prevent contributions from adding mechanically. Establish the shock's exclusion from unrelated productivity or policy news. Distinguish common shocks from bilateral spillovers and identify which dollar exposures create amplification rather than merely reflecting prior risk. Compare countries using prespecified initial states and horizon sets. Supply bounds when channel interventions are not separately observed. Do not equate principal components of global asset prices with an identified causal shock. Report whether the resulting decomposition is stable across crisis and ordinary conditions and across bank-dominated versus market-based financial systems.

\textbf{Additional formal definitions and scope (editorial).} A global shock is a scalar innovation in a specified safe-asset-demand equation, independent of the other innovations under stated exclusion assumptions; its unit-variance normalization must refer to the preintervention law. Dollar funding exposure is the quantity and maturity of dollar liabilities net of stated hedges, distinct from safe-asset demand. For a channel \(k\), define \(\tau_c^{(-k)}\) under an explicit equation-removal intervention and the channel contrast \(\tau_c-\tau_c^{(-k)}\). It need not equal a natural mediation effect or add to the total response across channels. Declare the timing of policy reactions and whether the country panel has interference through a network. Identification of a latent global factor up to sign/rotation does not identify this structural innovation or the channel contrasts.

For the identification part, declare an admissible class \(\Theta\) of structural models, including preferences, technologies, shock laws, measurement/sampling rules and any equilibrium selector. If \(O\) denotes the complete observable history and \(T(\theta)\) the target defined above, the identified set is \(\mathcal I_T=\{T(\theta):\theta\in\Theta,\ \mathcal L_\theta(O)=\mathcal L_{\rm obs}(O)\}\); point identification means this set is a singleton. A \(do(a)\) comparison replaces stated policy/technology equations while fixing the declared exogenous law and re-solving the remaining equations. These definitions do not assert that an identification theorem holds without further restrictions.
\subsection*{Variants and assumption changes}
\begin{enumerate}
\item \textbf{Observed exogenous global demand innovation (editorial).} Identify total country responses under a dated instrument and a fixed set of local policy rules.
\item \textbf{Channel interventions (editorial).} Introduce independently justified changes in dollar financing and nonbank transmission; estimate contrasts and their interaction under endogenous portfolios.
\end{enumerate}
\subsection*{References and source audit}
\begin{enumerate}
\item Official January 9, 2026 web version; locator: Interconnected World / Global drivers and 2026 international-propagation priority. Broad agenda; exact model editorial. URL: \url{https://www.bankofengland.co.uk/research/bank-of-england-agenda-for-research}.
\end{enumerate}
\begin{enumerate}\raggedright
\item\hypertarget{definitions:source:30007615:1}{} Bank of England. 2026. \emph{Bank of England Agenda for Research, 2025–2028}. Interconnected World / Global drivers and 2026 international-propagation priority.
\url{https://www.bankofengland.co.uk/research/bank-of-england-agenda-for-research}.
\end{enumerate}

\subsection*{Common conventions: Source mathematical conventions}

These conventions apply to each entry unless explicitly replaced.

\textbf{Models and identification.} A primitive model \(m\) specifies all probability laws, feasible choices, information, equilibrium and measurement rules used by its target. The admissible class \(\mathcal M\) is fixed independently of outcomes. If \(P_O(m)\) is its observable probability law and \(\tau(m)\) the target, its identified set at observable law \(P\) is
\[
 \mathcal I_\tau(P)=\{\tau(m):m\in\mathcal M,\ P_O(m)=P\}.
\]
Point identification means that this set is a singleton. An empty set signals incompatibility of data and assumptions. Coordinate bounds need not describe the joint set. Bounds are sharp only if every retained target is attainable or arbitrarily approachable by compatible models. Finite bounds require explicit restrictions on unobserved quantities; unrestricted confounding, hidden wealth, extrapolation or arbitrary equilibrium selection can yield uninformative sets.

\textbf{Causal interventions.} A potential outcome \(Y_i(p)\) is the outcome under a complete policy regime \(p\), including timing, financing, information and market spillovers. The target population and units are fixed before treatment. With interference, use the complete assignment vector or a justified exposure mapping; individual no-interference notation alone cannot identify an economy-wide intervention. A difference of mean potential outcomes does not require the cross-world joint distribution. Conditioning on post-treatment migration, survival, enrollment or employment changes the estimand and needs separate justification.

\textbf{Statistical statements.} An estimator, confidence set, forecast or test specifies a sampling law, model class and loss/coverage criterion. Uniform \((1-\alpha)\)-coverage means \(\inf_{m\in\mathcal M}\Pr_m\{\tau(m)\in C_n\}\geq1-\alpha\), or an explicitly asymptotic version. The whole parameter space trivially has coverage, so informativeness requires a stated width, risk or power objective. A forecasting improvement is relative to a named benchmark, information set and evaluation population.

\textbf{Equilibrium and optimization.} An equilibrium comprises feasible plans, optimality, beliefs where needed, budget constraints and all market/resource clearing. The primitives or a stated benchmark define each correspondence \(\mathcal E(p;m)\). Nonemptiness is assumed or proved, never inferred just from compact policy sets. A selected-equilibrium target uses a declared selection rule; a robust target quantifies over all permitted equilibria. Use suprema or infima unless attainment is established. An \(\varepsilon\)-optimal policy is feasible and within \(\varepsilon\) of the finite optimum value. Report an empty feasible set separately.

\textbf{Welfare and accounting.} Normative utility, interpersonal normalization, social weights, numeraire, discounting and the use of public receipts are primitives. Behavioral utility need not coincide with the welfare criterion. Household welfare includes ownership income and rents once; adding profits again to compensating variation generally double-counts them. A monetary equivalent is defined by an explicit transfer and price convention, with monotonicity and range conditions. Average effects, mechanism contributions, transfers and welfare losses are different targets. Infinite sums require convergence and dynamic feasible plans require terminal or no-Ponzi conditions.

\textbf{Source versions.} A working-paper date, revision date and journal publication year are distinguished. A DOI for a journal version is not automatically the DOI of an earlier manuscript. Locators refer to the stated version and printed pages where specified. A metadata-only check does not verify a claim about a full-text conclusion. Every entry's source subsection narrows the provenance claim accordingly.
% END ECONOMICS STATEMENT
\end{document}
